\documentclass[../main.tex]{subfiles}
\begin{document}
\section{Idempotent maps}\label{sec:idempotent}

% @title: Idempotent map
\begin{definition}\label[An idempotent self-map equals its own composite.]{def:idempotent}
A map $p\colon X\to X$ is \emph{idempotent} if $p\circ p=p$.
\end{definition}

% @title: The image is the fixed-point set
\begin{lemma}\label[The image of an idempotent map equals its fixed-point set.]{lem:image-fixed}
Let $p\colon X\to X$ be idempotent in the sense of \use{def:idempotent}.
Then
\[
  \im(p)=\{x\in X\mid p(x)=x\}.
\]
\end{lemma}
\begin{proof}[
  label={prf:image-fixed},
  proves={lem:image-fixed},
  description={Applies idempotence to image points and writes fixed points as their own images.}
]
If $x=p(y)$ for some $y\in X$, then $p(x)=p(p(y))=p(y)=x$ by \use{def:idempotent}.
Conversely, if $p(x)=x$, then $x=p(x)\in\im(p)$.
\end{proof}

\begin{remark}\label[The image and fixed-point sets are empty on the empty domain.]{rem:empty-set}
\mention{lem:image-fixed} also applies when $X$ is empty: both sets in its conclusion are empty.
\end{remark}
\end{document}
